
\subsection{A pressure-complete equation for the near-axis torque coefficient}
\label{sec:xlii-jet}

The leading Reynolds torque coefficient from Edition~XLI is a quadratic
functional of the Cartesian velocity jets, not a closed strain variable.
This section derives its exact equation from the \emph{full unforced}
Navier--Stokes system and provides an explicit, strictly positive
short-time amplification example. Its positive mechanism is
viscous redistribution; sustained terminal nonlinear amplification
remains unproved.

Let $u=(U,V,W)$ be a smooth divergence-free unforced solution on
$\mathbb R^3$. A subscript $0$ denotes evaluation at $x=y=0$.
Set $a=U_0$, $b=V_0$, $c=W_0$, and define
\[
A=-2(U_{xy})_0+(V_{xx})_0-(V_{yy})_0,\quad
B=(U_{xx})_0-(U_{yy})_0+2(V_{xy})_0,\quad
Q=b(W_x)_0-a(W_y)_0.
\]
The precise XLI quantity is
\begin{equation}\label{eq:xlii-lambda}
\boxed{\quad \Lambda=-\tfrac12(aA+bB+\partial_zQ).\quad}
\end{equation}
This equals $-4J_3-\partial_zK_2$, where
$\mathcal T=r^2\Lambda+O(r^4)$ is the angular-mean
Reynolds torque.

\begin{theorem}[Full projected Navier--Stokes jet evolution]
\label{thm:xlii-full}
For each transverse multiindex $\alpha$ and $i=1,2,3$, put
\begin{equation}\label{eq:xlii-H}
H_i^\alpha=\left[
\nu\Delta\partial_{xy}^\alpha u_i
-\partial_{xy}^\alpha\partial_i p
-\partial_{xy}^\alpha(u\cdot\nabla u_i)
+c\,\partial_z\partial_{xy}^\alpha u_i
\right]_0.
\end{equation}
Writing $D_0=\partial_t+c\partial_z$, one has
$D_0(\partial_{xy}^\alpha u_i)_0=H_i^\alpha$, and
\begin{equation}\label{eq:xlii-evolution}
\boxed{\begin{aligned}
D_0\Lambda=-\frac12\big[&
H_1^0 A+a(-2H_1^{xy}+H_2^{xx}-H_2^{yy})
+H_2^0 B+b(H_1^{xx}-H_1^{yy}+2H_2^{xy})\\
&+\partial_z\big(H_2^0(W_x)_0+bH_3^x
-H_1^0(W_y)_0-aH_3^y\big)
-(\partial_zc)(\partial_zQ)\big].
\end{aligned}}
\end{equation}
All pressure effects are retained via
\begin{equation}\label{eq:xlii-poisson}
-\Delta p=\partial_i\partial_j(u_iu_j),\qquad
p=\sum_{i,j}\mathcal R_i\mathcal R_j(u_iu_j),
\end{equation}
where $\mathcal R_i$ denotes a Riesz transform.
The equation generally needs fourth-order velocity
jets, third-order pressure derivatives, and nonlinear
products through third-order velocity jets. It is not
a closed autonomous equation in the second-order jets.
\end{theorem}
\begin{proof}
Differentiate the unforced physical momentum equation
$\partial_tu_i=\nu\Delta u_i-\partial_i p-u\cdot\nabla u_i$
before restricting to $x=y=0$, then add $c\partial_z$.
The product rule applied to \eqref{eq:xlii-lambda} gives
\eqref{eq:xlii-evolution}, noting the essential commutator
$D_0\partial_zQ=\partial_zD_0Q-c_z\partial_zQ$.
Taking divergence of the physical NSE gives
\eqref{eq:xlii-poisson}, including the nonlocal pressure
response. The derivative orders are explicit in
\eqref{eq:xlii-H}.
\end{proof}

\begin{theorem}[Strictly positive instantaneous amplification
in one globally regular unforced flow]
\label{thm:xlii-positive}
Fix any $\nu>0$, set $\varepsilon=\nu/1000$, and prescribe
the fully explicit nonzero Schwartz datum
\begin{equation}\label{eq:xlii-data}
\boxed{\quad u_0=\operatorname{curl}\left[
\varepsilon e^{-(x^2+y^2+z^2)}
\bigl(x-(x^2+y^2)y\bigr)e_z\right].\quad}
\end{equation}
It is exactly solenoidal and globally evolves as a unique
smooth, finite-energy, \emph{unforced} solution, by the
same Fourier small-data bound $q\le1/(100\sqrt\pi)<1/16$
used in Edition~XXXVII. Its exact initial near-axis
coefficient and time derivative satisfy
\begin{equation}\label{eq:xlii-amplification}
\boxed{\begin{aligned}
\Lambda(0,z)&=4\varepsilon^2e^{-2z^2},\\
\partial_t\Lambda(0,z)&=
16\nu\varepsilon^2(2z^2-7)e^{-2z^2},\\
\partial_t\Lambda(0,2)&=16\nu\varepsilon^2e^{-8}>0,\\
\partial_t\Lambda(0,0)&=-112\nu\varepsilon^2<0.
\end{aligned}}
\end{equation}
Consequently, for some $\delta>0$, not numerically
certified here, this one exact unforced solution has
\[
\boxed{\quad\partial_t\Lambda(t,2)
\ge 8\nu\varepsilon^2e^{-8}>0
\quad(0\le t\le\delta).\quad}
\]
The entire instantaneous nonlinear advection
\emph{including projected pressure} contributes
exactly zero to $\partial_t\Lambda(0,z)$ by
half-turn rotation parity. The positive rate at $z=2$
is viscous redistribution, not nonlinear strain growth.
\end{theorem}
\begin{proof}
Let $S=\operatorname{diag}(-1,-1,1)$ and
$\mathcal S v(x)=Sv(S^{-1}x)$.
The potential in \eqref{eq:xlii-data} is odd under
a half-turn, so $\mathscr S u_0=-u_0$.
The quadratic projected acceleration
$N(u_0)=-\mathbb P\nabla\cdot(u_0\otimes u_0)$
is even, $\mathscr S N(u_0)=N(u_0)$.
Since $\Lambda$ is an angular-average scalar,
it is a rotation-invariant quadratic jet functional.
Its polarization on an odd and an even velocity
must vanish: $d\Lambda_{u_0}[N(u_0)]=0$.
Thus the exact initial derivative arises only
from $\nu\Delta u_0$.

Independently calculate the heat evolution of the
Gaussian potential. Writing $s=1+4\nu t$ gives
\[
F_H=\varepsilon s^{-3/2}e^{-(x^2+y^2+z^2)/s}
\left[\frac{x}{s}-\frac{(x^2+y^2)y}{s^3}
-\frac{2(s-1)y}{s^2}\right].
\]
Its exact jet coefficient is
\[
\boxed{\quad\Lambda_H(t,z)=
4\varepsilon^2s^{-7}e^{-2z^2/s}.\quad}
\]
Differentiation at $t=0$ yields
\eqref{eq:xlii-amplification}. Time-continuity of
the smooth actual unforced solution provides
the strict positive bound for a sufficiently short interval.
\end{proof}

\paragraph{Release boundary.}
The results give one pressure-complete equation and one
strictly positive \emph{early-time} coefficient lower bound
from an exact unforced solution. They do not establish
a positive sustained nonlinear contribution, the required
source-compatible $\tau^{-h}$ terminal amplification, or
pointwise blowup. The Gaussian solution is globally regular.
No full Lean Navier--Stokes PDE formalization is claimed.
